\writeSectionFrame{\presentationTitle}{Das Problem}
\begin{frame}{Rückgängig machen}
	\begin{center}
 		\includegraphics[scale=0.3]{\BILDERPFAD/problem1.png}
 	\end{center}
\end{frame}

\note{
    Wir stellen uns einen Texteditor vor, in dem man Änderungen an einer	 Datei wieder rückgängig machen kann. Bevor Änderungen gespeichert werden, müssen wir also in irgendeiner Art und Weise einen Schnappschuss unserer aktuellen Version erstellen. Doch wie soll dieser Schnappschuss erstellt werden? Ganz offensichtlich müssen wir dafür Objekte kopieren.
}

\begin{frame}{Privater Zustand}
	\begin{center}
 		\includegraphics[scale=0.3]{\BILDERPFAD/problem2.png}
 	\end{center}
\end{frame}

\begin{frame}{Öffentlicher Zustand}
	\begin{center}
 		\includegraphics[scale=0.3]{\BILDERPFAD/problem3.png}
 	\end{center}
\end{frame}

\begin{frame}{Zusammenfassung}
	\begin{alertblock}{Privater Zustand}
 		Kopieren ist schwierig. 
 	\end{alertblock}
 	
 	\begin{alertblock}{Öffentlicher Zustand}
 		Unsicher
 	\end{alertblock}
\end{frame}

\begin{frame}{Wie speichern wir den Schnappschuss?}
	\begin{center}
 		\includegraphics[scale=0.3]{\BILDERPFAD/problem4.png}
 	\end{center}
\end{frame}

\note{
	Um den Schnappschuss zu speichern, werden wir irgendeine Art von Containerklasse	benötigen. Diese Klasse enthält alle Attribute der zu speichernden Klasse, also	eine ganze Menge Felder. Über die Methoden des Schnappschusscontainers sind dann wieder Zugriffe auf den Zustand des zu speichernden Objekts möglich. Dann hätten aber Änderungen in der zu speichernden Klasse wieder Auswirkungen auf die 
	Containerklasse. 
}